\documentclass[11pt]{article}
\usepackage[T1]{fontenc}
\usepackage{textcomp}
\usepackage[margin=0.75in]{geometry}
\usepackage{amsmath,amssymb,amsthm}
\usepackage{array,booktabs}
\usepackage{hyperref}
\setlength{\parindent}{0pt}
\newtheorem{theorem}{Theorem}
\theoremstyle{definition}
\newtheorem{definition}{Definition}
\title{Flow Proof Artifact: prop-03-proportional-cut-implies-parallel}
\author{Generated by \texttt{flow doc proof}}
\date{Flow Proof Book}
\begin{document}
\maketitle
If the sides of a triangle are cut proportionally, the join of the points of section is parallel to the remaining side.\\[0.5em]
\textbf{Source.} euclid — Elements, Book VI, Proposition 3\\[0.5em]
\bigskip
\noindent{\large \textbf{Derived fact 1.} \textit{If the sides of a triangle are cut proportionally, the join of the points of section is parallel to the remaining side} \hfill the Euclidean plane \cdot Euclid Book VI \cdot Proposition 3: a proportional cut implies a line parallel to the opposite side}\par
\smallskip\noindent\textit{Coordinate: the Euclidean plane \cdot Euclid Book VI \cdot Proposition 3: a proportional cut implies a line parallel to the opposite side}\par
\smallskip\noindent\textit{Source: euclid — Elements, Book VI, Proposition 3}\par
\smallskip\noindent\textit{Built on: proposition 2: a line parallel to a triangle side cuts the other sides proportionally, for Euclid Book VI on the Euclidean plane}\par
\medskip
\noindent\textbf{Goal.} If the sides of a triangle are cut proportionally, the join of the points of section is parallel to the remaining side\par
\noindent$\displaystyle \text{proportional sides implies parallel}$\par
\medskip
\noindent
\begin{tabular}{@{} >{\raggedright\arraybackslash}p{0.06\textwidth} >{\raggedright\arraybackslash}p{0.47\textwidth} >{\raggedleft\arraybackslash}p{0.41\textwidth} @{}}
\textbf{\#} & \textbf{Proof} & \textbf{Mathematics} \\
\midrule
\textbf{1.} & We prove that proposition 3: a proportional cut implies a line parallel to the opposite side for Euclid Book VI the Euclidean plane. & \\[0.45em]
\textbf{2.} & Let ABC = triangle. & \\[0.45em]
\textbf{3.} & Let D = point\_on\_AB. & \\[0.45em]
\textbf{4.} & Let E = point\_on\_AC. & \\[0.45em]
\textbf{5.} & From step 2, step 3, and step 4, we can deduce that DE parallel to BC. & $\displaystyle \text{DE parallel to BC}$ \\[0.45em]
\textbf{6.} & We invoke the derived fact governing Euclid Book VI the Euclidean plane: proposition 2: a line parallel to a triangle side cuts the other sides proportionally, for Euclid Book VI on the Euclidean plane. & \\[0.45em]
\textbf{7.} & From step 6, this implies proportional sides implies parallel. Hence proven. & $\displaystyle \text{proportional sides implies parallel}$ \\[0.45em]
\bottomrule
\end{tabular}
\label{thm:Geometry-Euclid Book VI-proposition_3_a_proportional_cut_implies_a_line_parallel_to_the_opposite_side}
\par\medskip\hrule
\end{document}
